\section{Introduction}\label{sec:introduction}

The Razak--Jacelon algebra \(\W\) is a simple, monotracial,
stably projectionless nuclear \(C^*\)-algebra
\cite{Jacelon}. Tensoring by \(\W\) removes the
\(K\)-theoretic part of the classification problem and makes the order
on the Cuntz semigroup visible through traces. The question is whether
the resulting trace data determine the algebra even when its ideal
structure is unrestricted. We prove that they do.

\subsection{The invariant and the main theorem}

For a \(C^*\)-algebra \(C\), let \(T(C)\) consist of all maps
\(\tau:C_+\to[0,\infty]\) satisfying
\[
 \tau(0)=0,\qquad
 \tau(a+b)=\tau(a)+\tau(b),\qquad
 \tau(ra)=r\tau(a)\quad(r>0),
\]
together with \(\tau(x^*x)=\tau(xx^*)\) and lower semicontinuity for
norm-convergent nets in \(C_+\). Arithmetic is extended nonnegative
arithmetic. We impose no density condition on the finite domain.
In particular, every closed two-sided ideal \(I\) contributes the weight
\begin{equation}\label{intro:ideal-weight}
 \tau_I(a)=
 \begin{cases}
  0,&a\in I_+,\\
  \infty,&a\in C_+\setminus I_+.
 \end{cases}
\end{equation}
The canonical topology is specified by
\begin{equation}\label{intro:trace-topology}
 \tau_i\longrightarrow\tau
 \quad\Longleftrightarrow\quad
 \limsup_i\tau_i((a-\eps)_+)
 \leq\tau(a)\leq\liminf_i\tau_i(a)
 \quad(a\in C_+,\ \eps>0).
\end{equation}
An isomorphism of these topological cones is a homeomorphism preserving
addition, the zero weight, and multiplication by every positive real
scalar.

\Needspace{3\baselineskip}
\begin{theorem}\label{thm:main}
Let \(A\) and \(B\) be separable nuclear complex \(C^*\)-algebras.
If \(T(A)\) and \(T(B)\) are isomorphic as topological cones, then
\[
 A\otimes\W\otimes\K\ \cong\ B\otimes\W\otimes\K,
\]
where \(\K=\K(\ell^2(\N))\) and all tensor products are spatial.
\end{theorem}

The conclusion is an actual \(C^*\)-algebra isomorphism.
The algebras may be nonunital, may have both finite and infinite
subquotients, and may have an arbitrary primitive ideal space.
Their only regularity assumptions are separability and nuclearity.
Moreover, the isomorphism can realize the prescribed cone map: after
transporting the given map to
\(F:T(B\otimes\W\otimes\K)\to T(A\otimes\W\otimes\K)\),
we obtain an isomorphism \(\Phi\) with
\(\tau\circ\Phi=F(\tau)\) for every extended trace \(\tau\).
This refinement is proved at the end of \Cref{pt:intertwining}.

The presence of \eqref{intro:ideal-weight} makes the invariant sensitive
to this ideal structure. For example,
\[
 I\subseteq J
 \quad\Longleftrightarrow\quad
 \tau_I+\tau_J=\tau_I.
\]
The finite and zero ideals of a general weight are also encoded by the
two limiting operations obtained by letting its positive scalar
multiple tend to zero and to infinity. These features will let us
localize each approximation on the ideal where its controlling trace
is finite.

\subsection{History and significance}

Razak's classification of simple inductive limits of stably
projectionless building blocks made the cone of traces, together with
its scale, an effective classification invariant
\cite[Theorems~1.1 and~2.1]{Razak}. Jacelon's construction singled out
the simple monotracial algebra $\W$ in this setting \cite{Jacelon}.
The existence, uniqueness, and approximate-intertwining organization
of those classification arguments also provides the broad method used
here, with new local estimates needed for arbitrary ideals.

The trace-cone question is due to Leonel Robert. It was recorded in
Santiago's 2012 conference abstract \cite{SantiagoWStable} and appears
as Problem~LXVIII in the survey of Schafhauser, Tikuisis, and White
\cite{STW}. Their formulation retains every extended trace, including
the zero/infinity ideal weights. Their May 2026 discussion describes
the known simple and traceless cases and the difficulty already
presented by one proper nonzero ideal. \Cref{thm:main} gives a positive
answer to that problem for the full cone and the unrestricted class of
separable nuclear algebras in its statement.

Elliott, Robert, and Santiago developed the compact topological cone
of extended lower-semicontinuous traces and its relation to
functionals on the Cuntz semigroup \cite{ERS}.
Robert subsequently identified tensoring with $\W$ with
realification of that semigroup \cite[Theorem~5.1.2]{Robert}.
Thus the cone in \eqref{intro:trace-topology} determines the Cuntz
semigroup of the stabilized algebra, including its ideals and
positive real scalar multiplication. The remaining task is to realize
this invariant by algebra maps and prove their uniqueness.

In the tracial simple case, the work of Elliott and Niu \cite{ElliottNiu}
and the published classification of Elliott, Gong, Lin, and Niu
\cite[Theorem~7.5]{EGLN} establish classification by scaled traces
under the stated $KK$-contractibility and finite-nuclear-dimension
hypotheses. Castillejos and Evington's nuclear-dimension theorem for
simple separable nuclear $\mathcal Z$-stable algebras supplies the
regularity bridge for the simple stabilized case
\cite[Theorem~A]{CastillejosEvington}. Nawata gave another proof that
$A\otimes\W\cong\W$ when $A$ is simple, separable, nuclear, has a
unique tracial state, and has no unbounded densely defined traces
\cite[Corollary~6.2(i)]{Nawata}. These simple-algebra hypotheses are
not imposed on $A$ and $B$ in \Cref{thm:main}.
As a separate regularity consequence, every stabilization
$A\otimes\W\otimes\K$ in \Cref{thm:main} has
nuclear dimension at most one, by \cite[Theorem~1.3]{OpenAINuclearDimension}
together with $\W\otimes\mathcal Z\cong\W$ from
\cite[Corollary~6.3]{Jacelon}.
The classification argument below does not use this nuclear-dimension
bound.

The traceless case has a different structure: when every extended
trace takes only zero and infinity, the trace cone reduces to ideal
information. R\o rdam's absorption criterion
\cite[Theorem~5.2]{RordamZ} and Kirchberg's ideal-lattice classification,
with a complete proof in \cite[Theorem~6.13]{GabeO2}, explain the known
traceless branch. The present argument must also retain finite trace
data and distinguish the regions where a weight is finite from those
where it is infinite. That distinction is necessary when finite and
infinite subquotients occur in the same algebra.

Restricted nonsimple precedents include Lin and Ng's classification
of certain essential extensions by $\W$, with a universal coefficient
theorem assumption, a simple quotient, and an invariant containing
$K$-theory and trace data \cite[Theorem~9.6]{LinNg}.
The general cone-only question allows arbitrary ideals.

The main analytic tools have separate roles. Ciuperca, Giordano, Ng,
and Niu give unital von Neumann algebra uniqueness from support-rank data
\cite{CGNN}; we combine this with trace caps and rational averaging to
obtain local rank estimates. The stable-uniqueness methods of
Dadarlat and Eilers \cite{DadarlatEilers} and the ideal-related
absorption and multiplier results of Gabe
\cite[Sections~11--13]{Gabe} remove the resulting errors in norm.
Those intermediate results apply here without pure infiniteness of
the auxiliary targets. In the existence argument, Connes's
hyperfiniteness theorem \cite{Connes}, nuclear completely positive
approximation \cite{ChoiEffros}, and measured-field methods
\cite{VaesWouters} provide the classical inputs. The fixed controllers,
lower ideal-support estimates, and fine-label correction that connect
these tools are proved below.

The strategy of lifting tracial data and then correcting a map through
an extension has a methodological precedent in Schafhauser's proof of
the Tikuisis--White--Winter theorem \cite{Schafhauser}. His setting
uses a faithful amenable trace and the universal coefficient theorem
(UCT); the present local error
ideals and support estimates address a different invariant, with no
UCT assumption. The use of quotient cone approximations and carriers
in prescribed ideals also has a precedent in
Bosa--Gabe--Sims--White \cite[Section~3]{BGSW}. Here the small UHF
fractions must satisfy both upper trace bounds and lower comparisons
into fixed positive cuts. The suspension-and-corner passage follows
the method in \cite[Remark~6.4]{GabeO2} and
\cite[proof of Theorem~14.1]{Gabe}; its trace normalization is proved
below for all extended weights.

\subsection{The construction}

Fix
\[
 P=A\otimes\W\otimes\K,\qquad
 Q=B\otimes\W\otimes\K.
\]
For a $C^*$-algebra $C$, the Cuntz semigroup $\Cu(C)$ consists of
Cuntz classes $[a]$ of positive elements in $C\otimes\K$, with addition
by orthogonal sum. The relation $a\precsim b$ means that
$v_n^*bv_n\to a$ in norm for some sequence $(v_n)$; it determines the
order on these classes. Two elements define the same class when each
is Cuntz subequivalent to the other. For a positive contraction $a$, its rank at an
extended trace is $d_\tau(a)=\lim_k\tau(a^{1/k})$, using unnormalized
matrix traces. This depends only on $[a]$.
Two obstacles distinguish this setting from one with a single
normalized trace. A positive element generating a stable algebra may
have infinite rank at every nonzero densely finite trace, so making
its rank a small scalar multiple gives no useful finite error bound.
Also, a weight can be finite only on a proper ideal. We must control
errors separately on these ideals and preserve enough lower support
to retain the infinite values outside them.

The given cone isomorphism yields compatible isomorphisms
\[
 F:T(Q)\longrightarrow T(P),\qquad
 G:\Cu(P)\longrightarrow\Cu(Q),
 \qquad
 d_\tau(G[a])=d_{F\tau}(a).
\]
We identify the primitive ideal spaces through this correspondence
and denote the common space by \(X\). For an open $U\subseteq X$,
write $P(U)$ and $Q(U)$ for the corresponding closed ideals.
We first construct maps into a sequence algebra, where errors may tend
to zero along the coordinates. Norm uniqueness will later allow their
coordinate maps to be conjugated into convergent sequences in $Q$.
Our intermediate maps take values in
\[
 D=Q_\infty=\ell^\infty(Q)/c_0(Q).
\]
We regard $Q$ as the constant sequences in $D$. If $\omega$ is a free
ultrafilter on $\N$ and $\tau_n\in T(Q)$, their regularized limit on
bounded positive sequences is
\[
 \rho([(x_n)])=\sup_{\eps>0}\lim_{n\to\omega}
                    \tau_n((x_n-\eps)_+),\qquad x_n\geq0.
\]
Section~\ref{sec:sequence-models} proves that this defines an extended
trace. We use $Y$ equal to a point, a half-open interval, or $\R$,
with a full-support Radon measure $m$, finite on the interval.
For $\nu\in T(P)$, write $\Fin(\nu)$ for the closed ideal generated by the positive
elements on which $\nu$ is finite. The product trace $m\otimes\nu$
is the integrated densely finite trace on $C_0(Y)\otimes\Fin(\nu)$,
extended to be infinity on positive elements outside this ideal.
On positive elementary tensors it has value $(\int f\,dm)\nu(a)$,
with $0\cdot\infty=0$.
A \emph{model} is a homomorphism
\(p:C_0(Y)\otimes P\to D\) satisfying
\begin{equation}\label{intro:model}
 \rho\circ p=m\otimes F(\rho|_Q)
\end{equation}
for every such regularized limit \(\rho\). The quantifier includes
nondensely finite weights. An ideal of $C_0(Y)\otimes P$ has a fine
support, an open subset of $Y\times X$; its projection to $X$ is the
coarse support used to index local estimates.

The main difficulty is to preserve both the finite and the infinite
parts of \eqref{intro:model}. Fix a regularized limit $\rho$, put
$\sigma=\rho|_Q$, and write $Q(V)=\Fin(\sigma)$. On
$C_0(Y)\otimes P(V)$, controlled trace moments will identify the
pullback of $\rho$. This alone does not prevent the pullback from
becoming finite on a larger ideal. A separate lower ideal comparison
rules that out. Once the finite ideal is exactly
$C_0(Y)\otimes P(V)$, both sides of \eqref{intro:model} are infinite
on every positive element outside it. This two-part recovery is the
reason for the following constructions.

\paragraph{Local error ideals and stabilizers.}
For each open \(U\subseteq X\), we define an ideal \(J(U)\subseteq D\)
using uniform bounds on the ranks of positive cuts.
The error at each cut is a scalar tending to zero times the rank of
a fixed positive controller $c$ satisfying
$c\precsim(b-\delta)_+$ for some $b\in Q(U)_+$ and $\delta>0$.
Such errors are invisible to every regularized weight that is densely
finite on \(Q(U)\).
We construct small homomorphisms whose values belong to these ideals
and dominate any prescribed countable family of smaller errors.
Compact containment of open ideals supplies a common controller
when different approximation stages are combined.
The two versions in Section~\ref{sm:section} have different inputs:
one takes a small fraction of an existing model; the other builds a
small cone homomorphism before any model is available.

\paragraph{Tracial approximation with lower ideal support.}
Convex separation and von Neumann algebra approximation construct
completely positive cone maps with the required first and second
trace moments. The construction simultaneously imposes norm tests
whose positive averages force lower rank comparisons.
The comparison constant is independent of the averaging
multiplicities. This lower signal determines the part of the cone
on which a weight must remain infinite.
Factor representations and their measurable assembly are used only
through cut tests controlled by fixed finite-rank constants.

\paragraph{Correction with all ideal labels retained.}
A small stabilizer places the multiplicative errors in a
sigma-unital hereditary target. Module Stinespring dilation and
absorption then correct the cone map.
The support estimate is imposed on the entire projected ideal of
each input. It therefore survives scaling of the cone variable and
makes every dilation weakly equivariant for the finer primitive-space
labels. Compact corrections can be compressed to finitely many terms
of the original diagonal repeat, with an explicit tail estimate.
A second small homomorphism then dominates the correction itself.
Adding it back retains the lower ideal comparison, so the finite and
infinite parts combine to give the full model equation in
Section~\ref{cc:section}.

\paragraph{Uniqueness and passage to actual maps.}
For two models with the same data, tracial alignment first makes their
difference locally small. A stabilizer drawn from a spare model then
allows ideal-related stable uniqueness to remove that difference in
norm. The spare is supplied by splitting each model into two UHF
halves and replacing the halves successively. This argument, in
Section~\ref{sec:uniqueness}, assumes that the two models are given;
it does not depend on the later existence theorem.
Once corrected cone models exist, an
exhaustion of the line gives models with Lebesgue data.
Translation covariance and a full crossed-product projection give
a model for \(P\) itself. A finite positive partition identity computes
the trace of this corner, including infinite values.
Uniqueness for every pair of subsequences gives eventual conjugacy
between any two sufficiently late coordinate maps. We use these
conjugacies to make a subsequence Cauchy in norm on each input; its
limit is an actual homomorphism $P\to Q$. Repeating the construction
in the other direction gives a map $Q\to P$. Their compositions have the
identity data, so approximate intertwining proves \Cref{thm:main}.

Figure~\ref{fig:proof-roadmap} displays these dependencies. In
particular, the uniqueness theorem is conditional on a pair of models;
it is available when the existence construction later supplies them.

\begin{figure}[tbp]
\centering
\begin{tikzpicture}[
  stage/.style={draw=MidnightBlue!65,fill=MidnightBlue!3,
    rounded corners=1pt,align=center,font=\small,
    text width=54mm,minimum height=11mm,inner sep=3pt},
  arrow/.style={-{Latex[length=2mm]},draw=MidnightBlue!80,thick}
]
\node[stage,text width=67mm] (data) at (0,0)
  {Full trace cone\\Realification and ideal transport (Section~\ref{sec:structure})};
\node[stage,text width=67mm] (errors) at (0,-1.6)
  {Fixed-cut comparisons and local error ideals\\(Section~\ref{sec:sequence-models})};
\node[stage] (small) at (-3.5,-3.3)
  {Small stabilizers, with or without\\a given model (Section~\ref{sm:section})};
\node[stage] (rough) at (3.5,-3.3)
  {Tracial approximation on the cone\\(Section~\ref{tc:section})};
\node[stage] (unique) at (-3.5,-5.0)
  {Norm uniqueness for given models\\(Section~\ref{sec:uniqueness})};
\node[stage] (cone) at (3.5,-5.0)
  {Fine-label correction and\\retained lower support (Section~\ref{cc:section})};
\node[stage,text width=67mm] (point) at (0,-6.8)
  {Covariant line model and point model\\(Section~\ref{pt:section})};
\node[stage,text width=67mm] (actual) at (0,-8.4)
  {Actual homomorphisms and intertwining\\(Section~\ref{pt:section})};
\draw[arrow] (data) -- (errors);
\draw[arrow] (errors.south west) -- (small.north);
\draw[arrow] (errors.south east) -- (rough.north);
\draw[arrow] (small) -- (unique);
\draw[arrow] (rough) -- (cone);
\draw[arrow] (small.east) -- (0,-3.3) -- (0,-5.0) -- (cone.west);
\draw[arrow] (unique.south) -- (point.north west);
\draw[arrow] (cone.south) -- (point.north east);
\draw[arrow] (point) -- (actual);
\draw[arrow,dashed] (unique.west) -- (-7.0,-5.0) -- (-7.0,-8.4) -- (actual.west);
\end{tikzpicture}
\caption{Principal model-construction dependencies. The existence and
uniqueness arguments meet in the passage from sequence models to actual
maps. Shared auxiliary inputs are not drawn separately: the trace-cap
and barycenter lemmas in Section~\ref{sec:uniqueness} are used in
Section~\ref{tc:section}; the absorption and hereditary-support lemmas
of Section~\ref{sec:uniqueness} are used in Section~\ref{cc:section}.
Norm uniqueness is a statement about
models when they are given; its proof does not assume the cone-existence
theorem. The dashed arrow records its further use in actual-map
extraction and the final two-sided intertwining.}
\label{fig:proof-roadmap}
\end{figure}


\subsection{Organization and conventions}

\Cref{st:realification,st:ideals} collect the structural input.
The sequence comparisons and local diagonalization are proved in
\Cref{seq:sandwich,seq:diagonal}. The stabilizer construction is
\Cref{sm:stabilizers}, and norm uniqueness is
\Cref{un:uniqueness}. The two existence stages are
\Cref{tc:rough-model,cc:cone-existence}; the final passage to an
isomorphism is \Cref{pt:intertwining}.

All algebras and Hilbert spaces are complex. A completely positive
contractive map is abbreviated c.p.c. Traces on finite matrix
amplifications use the unnormalized matrix trace.
The universal UHF algebra is denoted by \(\UHF\), to distinguish it
from the target algebra \(Q\). Rational diagonal averaging uses the
normalized UHF trace only in the averaging direction.
The notation \(V\Subset U\) denotes compact containment in the
lattice of open subsets of \(X\), without a Hausdorff assumption.
